%% ma: L11-B17-0017
%% lop: 11
%% chuong: 5
%% bai: 17
%% dang: 11.17.3
%% loai: TL
%% muc: VD
%% dap_an: "-"
%% nguon: "tai-lieu-hoc-tap-mon-toan-lop-11-theo-chuong-trinh-giao-duc-pho-thong-2018.pdf (Lớp Toán Thầy Duy, Phạm Lê Duy, THPT Chu Văn An), Chương 3 Bài 3, Dạng 3, Ví dụ 17 (tr.393-398)"
%% kien_thuc: "Lớp 11 Bài 17 (hàm số liên tục)"
%% trang_thai: chinh-thuc
%% ghi_chu: "Đã duyệt 03/11/2026."
\begin{ex}
Chứng minh rằng phương trình $x^3-2x^2+x-1-\dfrac2x+\dfrac m{x+3}=0$ luôn có nghiệm với mọi $m$.
\loigiai{\sloppy Điều kiện $x\ne0,\ x\ne-3$. Nhân hai vế với $x(x+3)\ne0$ ta được phương trình tương đương $g(x)=x(x+3)\left(x^3-2x^2+x-1\right)-2(x+3)+mx=0$. $g$ là đa thức bậc $5$ hệ số cao nhất dương, liên tục trên $\mathbb R$, $g(0)=-6<0$ và $g(t)\to+\infty$ khi $t\to+\infty$, nên tồn tại $b>0$ với $g(b)>0$. Khi đó $g(0)g(b)<0$ nên $g(x)=0$ có nghiệm $x_0\in(0;b)$; vì $x_0>0$ nên $x_0\ne0,\ x_0\ne-3$ và là nghiệm của phương trình đã cho.}
\end{ex}
